\documentclass{article}
% Top margin, right margin, left margin, bottom margin, footnote skip
\usepackage[utf8]{inputenc}
\usepackage{biblatex}
\addbibresource{./references.bib}
% linktocpage shall be added to snippets.
\usepackage{hyperref,theoremref}
\hypersetup{
	colorlinks, 
	linkcolor={red!40!black}, 
	citecolor={blue!50!black},
	urlcolor={blue!80!black},
	linktocpage % Link table of content to the page instead of the title
}

\usepackage{blindtext}
\usepackage{titlesec}
\usepackage{amsthm}
\usepackage{thmtools}
\usepackage{amsmath}
\usepackage{amssymb}
\usepackage{graphicx}
\usepackage{titlesec}
\usepackage{xcolor}
\usepackage{multicol}
\usepackage{hyperref}
\usepackage{import}


\newtheorem{theorem}{Theorema}[section]
\newtheorem{lemma}[theorem]{Lemma}
\newtheorem{corollary}{Corollarium}[section]
\newtheorem{proposition}{Propositio}[theorem]
\theoremstyle{definition}
\newtheorem{definition}{Definitio}[section]

\theoremstyle{definition}
\newtheorem{axiom}{Axioma}[section]

\theoremstyle{remark}
\newtheorem{remark}{Observatio}[section]
\newtheorem{hypothesis}{Coniectura}[section]
\newtheorem{example}{Exampli Gratia}[section]

% Proof Environments
\newcommand{\thm}[2]{\begin{theorem}[#1]{}#2\end{theorem}}

%TODO mayby proof environment shall have more margin
\renewenvironment{proof}{\vspace{0.4cm}\noindent\small{\emph{Demonstratio.}}}{\qed\vspace{0.4cm}}
% \renewenvironment{proof}{{\bfseries\emph{Demonstratio.}}}{\qed}
\renewcommand\qedsymbol{Q.E.D.}
% \renewcommand{\chaptername}{Caput}
% \renewcommand{\contentsname}{Index Capitum} % Index Capitum 
\renewcommand{\emph}[1]{\textbf{\textit{#1}}}
\renewcommand{\ker}[1]{\operatorname{Ker}{#1}}

%\DeclareMathOperator{\ker}{Ker}

% New Commands
\newcommand{\bb}[1]{\mathbb{#1}} %TODO add this line to nvim snippets
\newcommand{\orb}[2]{\text{Orb}_{#1}({#2})}
\newcommand{\stab}[2]{\text{Stab}_{#1}({#2})}
\newcommand{\im}[1]{\text{im}{\ #1}}
\newcommand{\se}[2]{\text{send}_{#1}({#2})}

\title{Ars Amatoriae}
\author{Publius Ovidius Naso} 
\date{\today}

\begin{document}
% \tableofcontents

Dear Awards Panel of China Oxford Scholarship Fund,
\\

My proposed DPhil research focuses on quantum computation. 
Computation, regardless of whether it is performed by the ancient abacus or the modern electronic computers, at its fundamental level is a physical process. 
While the abacus and electronic computers are based on classical physics, quantum computers perform computation by exploiting principles of quantum mechanics. 
Quantum mechanics permits phenomena that seem impossible in classical perspective, such as the superposition of the dead and alive states of Schrödinger’s cat. 
It is precisely those bizarre phenomena that allow us to break the classical constraints and perform computations that were previously thought to be impossible. 
Theoretical studies have shown that some quantum algorithms, such as Shor’s algorithm for factoring, are exponentially faster than their best classical equivalents. 
This ability of quantum computers to outperform classical ones is called quantum supremacy.

While the theory of quantum computations is promising, its physical realisation is extremely challenging.
At the moment, quantum computer can only handle dozens of qubits, which is the quantum mechanical equivalent of a bit. 
This number is dwarfed by the trillions of bits possessed by modern electronic computers. 
Nevertheless, quantum supremacy in specific tasks has been achieved even with such a limited system.
In 2024, Google reported its Willow quantum computer with 105 qubits completed a random circuit sampling task in five minutes that would take today's fastest supercomputers 10 septillion ($10^25$) years \cite{google}. 


Because of these promising results, many people, including myself, believe that quantum computers will fundamentally change our ways of computation and revolutionise our society like electronic computers have done. 
Many corporations, including Huawei, Baidu, Google, and Microsoft have formed their quantum computation divisions. 
Billions of investments and hundreds of startup companies are following them closely. 
Policy makers are introducing initiatives to sponsor research in quantum technologies. 
The UK has established the National Quantum Technology Programme, aiming to foster a quantum enabled economy. 
Very importantly, in China’s recent 14th National Committee of the Chinese People's Political Consultative Conference in March 2026, Premier Li Qiang stated that quantum technology is one of the new technologies of future and quality productive force. 
Through my DPhil study in quantum computation, I hope to make meaningful contributions to this exciting field of research. 
After completing my study, I plan to follow Premier Li’s call by returning to China, founding a quantum computation company, and applying my skills and knowledge in quantum computation to make positive impacts on our society. 

On February 18 2026, I was honoured to receive an offer from the University of Oxford Department of Computer Science to study DPhil under the supervision of Professor Sergii Strelchuk and Dr Sathyawageeswar Subramanian. 
Unfortunately, I have not secured any scholarship or funding. I would be extremely grateful if COSF may support my study, and, ultimately, my goal to bring positive impacts to our society through research in quantum information.



\hfill Yuhao Han \hspace{1.6em}

\hfill \today

\printbibliography
\end{document}
